% =====================================================================
%  subsec-uniqueness.tex
%  Uniqueness of the K-ary lax monoidal structure (Jon's task 2).
%
%  DESTINATION: sourcehut/algebraic-elaboration/main.tex
%  PLACEMENT:   as a SUBSECTION at the end of
%                 \section{The K-ary lax monoidal structure of a relative monad}
%                 (label sec:metalanguage-new), i.e. immediately after
%                 Definition~\ref{def:k-ary-relative} (and after the explanatory
%                 remarks of sec-metalanguage-new.tex, if those are added), just
%                 before \section{Modular elaborators in the monadic metalanguage}.
%
%  PREAMBLE DEPENDENCIES (all already present in the paper):
%     amsmath (\varprojlim); amssymb (\twoheadrightarrow); tikz-cd;
%     the theorem environments proposition / lemma / theorem / remark;
%     the macro \PMC (= \mathbf{L}_\Sigma).
%  CROSS-REFERENCES: def:k-ary-relative and con:lax-monoidal-compact,
%     both already in the paper.
%
%  BIB: the finite-uniqueness Proposition cites the two entries below;
%       add them to refs.bib (verify the MSFP page range):
%     @article{kock:1972,
%       author  = {Kock, Anders},
%       title   = {Strong functors and monoidal monads},
%       journal = {Archiv der Mathematik}, volume = {23}, number = {1},
%       pages   = {113--120}, year = {1972}, doi = {10.1007/BF01304852}}
%     @inproceedings{mcdermott-uustalu:2022,
%       author    = {McDermott, Dylan and Uustalu, Tarmo},
%       title     = {What Makes a Strong Monad?},
%       booktitle = {Proc.\ 9th Wksh.\ on Mathematically Structured Functional
%                    Programming (MSFP 2022)}, series = {EPTCS},
%       volume    = {360}, pages = {113--133}, year = {2022},
%       doi       = {10.4204/EPTCS.360.6}}
% =====================================================================

\subsection{Uniqueness of the $\mathcal{K}$-ary structure}
\label{sec:k-ary-uniqueness}

The structure of Definition~\ref{def:k-ary-relative}, when it exists, is unique. The value it produces is pinned down by naturality and the finite structure; its only genuinely infinitary datum---the definedness it imposes, which lives on $T\mathbf 1$---is forced to be an infinitary meet. Throughout this subsection we assume in addition that $T$ is \emph{commutative} and that $\mathcal{K}$ is closed under finite products and under passage to subsets of its members. The argument uses only Unitality, Associativity and Re-indexing (axioms (1)--(3) of Definition~\ref{def:k-ary-relative}), naturality of $\phi$, and the finite input below; the unit- and bind-compatibility axioms (4)--(5) are never invoked.

\begin{proposition}[Finite uniqueness]\label{prop:finite-uniqueness}
For a commutative relative monad over a base on which tensorial strengths are unique---in particular over $\mathbf{Set}$---the structure $\phi_{K,X}$ is uniquely determined at every \emph{finite} arity $K$. This is the relative form of the strength-uniqueness results of \citet{kock:1972} and \citet{mcdermott-uustalu:2022}.
\end{proposition}

Write $\mathbf 1$ also for the constant family at the terminal object of $\mathbf{Set}$. For an arity $S\in\mathcal{K}$, evaluating $\phi$ there yields (since $\prod_{s\in S}\mathbf 1\cong\mathbf 1$) a morphism of $\mathcal{C}$
\[
  d_S\;:=\;\phi_{S,\mathbf 1}\;\colon\;(T\mathbf 1)^{S}\longrightarrow T\mathbf 1,\qquad m:=d_{\mathbf 2},
\]
the \emph{combinator} at arity $S$. We argue with generalized elements: a \emph{point at stage $A$} is a morphism $A\to T\mathbf 1$ of $\mathcal{C}$, and a point $f\colon A\to(T\mathbf 1)^{S}$ has components $f_s:=\mathrm{pr}_s\circ f$.

\begin{lemma}\label{lem:terminal-meet}
The operation $m$ is commutative, associative and idempotent, so $p\sqsubseteq q:\Longleftrightarrow m\circ\langle p,q\rangle=p$ is a partial order on each hom-set $\mathcal{C}(A,T\mathbf 1)$ for which $m$ computes binary meets. For every $S\in\mathcal{K}$ the combinator $d_S$ is a lower bound \textnormal{(}$d_S(f)\sqsubseteq f_s$ for each $s$\textnormal{)}, fixes constant families \textnormal{(}$d_S\circ\Delta_S=\mathrm{id}$, with $\Delta_S$ the diagonal\textnormal{)}, and is monotone.
\end{lemma}

\begin{proof}
Commutativity, associativity and idempotence of $m$ are Re-indexing along the transposition $\mathbf 2\to\mathbf 2$, Associativity on the two bracketings of $\mathbf 3$, and Re-indexing along the fold $\nabla\colon\mathbf 2\twoheadrightarrow\mathbf 1$ (there the reindexing map is the diagonal and $\phi_{\mathbf 1,\mathbf 1}=\mathrm{id}$ by Unitality, so the square reads $\mathrm{id}=m\circ\Delta$); these make $\sqsubseteq$ a partial order with $m$ the meet. \emph{Lower bound:} Associativity on $S=\{s\}\sqcup(S\setminus\{s\})$ gives $d_S=m\circ(\mathrm{id}\times d_{S\setminus\{s\}})$, whence $d_S(f)\sqsubseteq f_s$. \emph{Constants:} Re-indexing along $S\twoheadrightarrow\mathbf 1$, as for idempotence. \emph{Monotonicity:} computing $d_{S\times\mathbf 2}$ by Associativity for its two factorizations $\coprod_{s\in S}\mathbf 2=\coprod_{i\in\mathbf 2}S$ (matched by Re-indexing along $S\times\mathbf 2\cong\mathbf 2\times S$) gives the interchange law $m\circ(d_S\times d_S)=d_S\circ m^{S}$; if $f\sqsubseteq g$ pointwise then $m\langle d_Sf,d_Sg\rangle=d_S\circ m^{S}\circ\langle f,g\rangle=d_S(f)$.
\end{proof}

\begin{lemma}[Rigidity of the combinator]\label{lem:combinator-rigidity}
For any structure, $d_S(f)$ is the greatest lower bound $\bigwedge_{s\in S}f_s$ in the order $\sqsubseteq$. Consequently any two structures with the same finite part have the same combinators: $d_S=d'_S$ for every $S\in\mathcal{K}$.
\end{lemma}

\begin{proof}
$d_S(f)$ is a lower bound (Lemma~\ref{lem:terminal-meet}) and it is the greatest: if $g\sqsubseteq f_s$ for all $s$---equivalently $\Delta_S\circ g\sqsubseteq f$ pointwise---then $g=d_S(\Delta_S\circ g)\sqsubseteq d_S(f)$ by monotonicity and constant-fixing. If $\phi,\phi'$ share their finite part then $m=d_{\mathbf 2}=d'_{\mathbf 2}$ (arity $\mathbf 2$ is finite), so they induce the same order $\sqsubseteq$; both $d_S(f)$ and $d'_S(f)$ are then greatest lower bounds of the same family, and a greatest lower bound is unique.
\end{proof}

\begin{theorem}[Uniqueness of the $\mathcal{K}$-ary structure]\label{thm:k-ary-uniqueness}
Let $K\in\mathcal{K}$ be an arity such that $T$ preserves, for every family $X\colon K\to\mathbf{Set}$, the cofiltered limit presenting the infinite product as the limit of its finite subproducts,
\[
  \prod_{k\in K}X_k\;=\;\varprojlim_{F\in\mathcal{P}_{\mathrm{fin}}(K)}\ \prod_{j\in F}X_j
\]
\textnormal{(}transition maps the projections\textnormal{)}. Then any two $\mathcal{K}$-ary lax monoidal structures on $T$ agree at $K$; in particular the structure on $T$, if it exists, is unique.
\end{theorem}

\begin{proof}
Let $\phi,\phi'$ be two structures. They agree at finite arities (Proposition~\ref{prop:finite-uniqueness}), so $d_S=d'_S$ for all $S$ (Lemma~\ref{lem:combinator-rigidity}). Fix $X$. Since $T$ preserves the displayed limit, the cone $\{T(\pi_F)\}_{F\in\mathcal{P}_{\mathrm{fin}}(K)}$ is a limit cone, hence jointly monic; it suffices to show $T(\pi_F)\circ\phi_{K,X}=T(\pi_F)\circ\phi'_{K,X}$ for each finite $F\subseteq K$. Let $X^{F}\colon K\to\mathbf{Set}$ agree with $X$ on $F$ and equal $\mathbf 1$ off $F$, and let $h\colon X\to X^{F}$ be the identity on $F$ and the unique map to $\mathbf 1$ off $F$; then $\prod_K h=\pi_F$ under $\prod_K X^{F}\cong\prod_F X$, so naturality of $\phi$ gives
\[
  T(\pi_F)\circ\phi_{K,X}\;=\;\phi_{K,X^{F}}\circ\textstyle\prod_k T(h_k).
\]
By Associativity on the partition $K=F\sqcup(K\setminus F)$,
\[
  \phi_{K,X^{F}}\;=\;\phi_{\mathbf 2,W}\circ\bigl(\phi_{F,X|_F}\times d_{K\setminus F}\bigr),\qquad W=\bigl(\textstyle\prod_F X,\ \mathbf 1\bigr),
\]
up to the canonical isomorphisms. The finite factors $\phi_{F,X|_F}$ and $\phi_{\mathbf 2,W}$ are equal for $\phi$ and $\phi'$ by Proposition~\ref{prop:finite-uniqueness}, and $d_{K\setminus F}=d'_{K\setminus F}$ by Lemma~\ref{lem:combinator-rigidity}; joint monicity gives $\phi_{K,X}=\phi'_{K,X}$.
\end{proof}

\begin{remark}
The proof uses only Unitality, Associativity and Re-indexing, naturality, finite uniqueness, and preservation of the cofiltered limits; the unit- and bind-compatibility axioms and any order-enrichment of $\mathcal{C}$ are not used. Uniqueness is thus a property of the underlying commutative $\mathcal{K}$-ary lax monoidal \emph{functor}, for any finite-product-preserving root $J\colon\mathbf{Set}\to\mathcal{C}$. In the concrete instance of Construction~\ref{con:lax-monoidal-compact}---$\mathcal{C}=\mathbf{Pos}$ (or the ambient topos), $J$ the discrete embedding, $T=\PMC$ the partiality monad---one has $T\mathbf 1\cong\Sigma$, the operation $m$ is binary conjunction, and $d_S=\forall_S$ is the $\mathbf{scope}$ combinator $\boldsymbol{\vartheta}_{S,\mathbf 1}$. Re-indexing holds only along surjections: along a projection one already has $T(\pi_1)\circ\phi_{\mathbf 2}\neq\pi_1$ for the lift monad, since $\phi_{\mathbf 2}$ has conjoined the definedness that $\pi_1$ discards.
\end{remark}
